%!TEX root = ../main.tex

\chapter{总结和展望}

\section{技术总结}

本文聚焦于文本属性图的理解与推理任务，旨在解决现有图结构大模型在结构编码信息损失、
跨模态表征空间割裂以及复杂推理过程不可追溯等核心挑战 。
本文设计实现了一种面向节点分类和图推理的图结构大模型关键技术框架，
实现了基于欧拉路径的可逆图序列化编码与多层次跨模态对齐机制，
实现并测评了基于可验证奖励的强化学习推理策略优化方法，
探索了涵盖学术引文、电商共购及合成图论任务的混合数据集，
通过结构保真表征与结合思维链冷启动的强化学习微调提升了模型在节点分类与复杂图推理任务上的性能，
并通过金融反欺诈原型系统验证了改进方案的有效性 。最后总结梳理文本所做的工作与成果：

1）本文设计实现了一种结合欧拉路径图表征和跨模态对齐的节点分类方案，
结合无损图序列化编码与多层次对齐损失函数。
该方案针对传统邻域采样导致的信息丢失问题，引入基于欧拉路径的可逆编码机制，
实现了子图拓扑的近似无损表征；同时，通过融合全局InfoNCE对比、
子结构-短语级局部匹配及生成式重建损失，
有效改善了图拓扑特征与文本语义特征之间的异构鸿沟，
提升了图模型的结构感知能力。

2）本文提出了一种结合强化学习微调的图推理增强技术，
构建了思维链监督冷启动与基于GRPO的可验证强化学习两阶段优化框架。
针对图推理任务缺乏高质量标注且泛化性差的问题，
本文利用合成图论任务构建大规模思维链数据进行冷启动微调，
进而设计了包含答案正确性、拓扑一致性与约束满足度的可验证奖励函数。
通过组式策略优化算法，在无需人工偏好标注的前提下，
引导模型生成逻辑自洽、步骤可核验的推理链，
在Erdős及GraphArena等评测中提升了模型处理复杂图逻辑的准确率与鲁棒性 。


3）本文设计与实现了一个融合节点分类与图推理的金融反欺诈原型系统，
验证了所提技术在垂直领域动态异构图风险识别场景中的应用价值。
系统基于预训练图结构大模型，
构建了从全量风险感知到深度归因推理的级联式分析链路。
通过适配金融交易流的欧拉路径编码与交互式可视化设计，
该系统不仅实现了风险节点判别，
还能输出包含关键交易路径与团伙特征的归因证据链，
为传统风控模型可解释性不足问题提供了解决方案。

\section{未来工作展望}

本文虽然在图结构大模型的表征优化、跨模态对齐及推理增强方面取得了一定成果，并成功应用于金融反欺诈场景，但面对日益复杂的图数据形态与更深层次的认知推理需求，仍存在若干值得进一步探索的研究方向：

1）面向超大规模图的长程依赖建模与效率优化。本文提出的基于欧拉路径的图序列化方法虽然实现了结构的无损编码，
但其序列长度随图规模呈线性增长，受限于大语言模型的上下文窗口。
未来工作可以尝试层次化图压缩编码或线性注意力机制，在保持拓扑保真度的同时降低序列长度。
此外，针对十亿节点级别的超大规模动态图，可以引入图检索增强生成技术，实现按需检索关键子结构。

2）增强复杂逻辑推理的稳健性与自适应性。在第四章的实验中，尽管引入了强化学习微调，
但模型在处理复杂图论任务及 NP-hard 类组合优化问题时仍有提升空间。
未来研究可借鉴思维树或蒙特卡洛树搜索与大模型结合的思路，赋予模型回溯与前瞻能力。
此外，针对当前强化学习中奖励权重与 KL 正则系数依赖网格搜索的问题，
未来计划引入基于自适应拉格朗日乘子的动态调节策略，通过在训练过程中自动优化约束边界，
取代人工调优，从而提升模型在不同复杂程度任务下的训练稳定性与泛化能力。

3）隐私保护下的图结构大模型安全建模。鉴于金融、医疗等实际场景对数据隐私的高度敏感性
，未来的研究需进一步权衡模型性能与数据安全。一方面，计划引入差分隐私或同态加密技术，
在节点文本脱敏的基础上，实现图结构与特征的深度隐私化处理；另一方面，将探索联邦图学习
框架下的跨机构联合建模方案，
在不交换原始数据的前提下实现图结构大模型的分布式协同训练，解决“数据孤岛”问题。

4）拓展多模态异构图数据的统一建模。目前的跨模态对齐主要聚焦于文本-图结构对。
未来计划将图结构大模型的输入端扩展为任意模态，设计通用的多模态投影接口，
使模型能够联合理解视觉、文本、时序流与拓扑结构中的关联信息，构建更加立体的图智能感知系统。

尽管如此，本文所做的工作依然是领域内的积极尝试。从结构序列化到语义对齐，
再到推理增强的研究路径，不仅验证了图—语言模型统一框架的可行性，
也将为后续面向工业大规模应用的研究提供参考。